\documentclass[10pt,a4paper]{article}
\usepackage[utf8]{inputenc}
\usepackage{amsmath}
\usepackage{amsfonts}
\usepackage{amssymb}
\usepackage{german}
\usepackage{graphicx}
\parindent 0pt
\pagestyle{headings}


\title{
	\includegraphics[width=10cm]{logo.png} \\
	\vspace{1cm}
	Pflichtenheft
}
\author{SEP - ITS - Team \\ Max Binder, Florian Büchner, Martin Freund, \\Florian Lorenz,
											Andreas Poxrucker, Andreas Wilhelm}
\begin{document}
	\maketitle
	\newpage
	\tableofcontents	
	\newpage
	
	\section{Zielbestimmung}
		Das Ziel des Projekts ist die automatische Erkundung eines unbekannten, rechtwinkligen Spielfelds durch bis zu
		sechs e-pucks. Die Erkundung soll durch Zusammenarbeit der Roboter möglichst effizient erfolgen.
		
		Das bereits erkundete Gebiet sowie die aktuellen Positionen der e-pucks werden kontinuierlich auf einem
		Andorid Smartphone dargestellt.
		
		Außerdem kann mit dem Smartphone einer der teilnehmenden e-puck Roboter ausgewählt und halbmanuell gesteuert werden.
		\subsection{Musskriterien}
			\begin{itemize}
				\item e-puck Roboter
				\begin{itemize}
					\item Variable Anzahl von Robotern
						\\ \textsl{Die Erkundung der Spielfläche kann durch bis zu sechs e-puck Roboter erfolgen.}
					\item Kommunikation über ein Bluetooth-Netzwerk
						\\ \textsl{Die Kommunikation der e-pucks untereinander erfolgt über ein Bluetooth-Netzwerk.
						   Dieses wird von den e-pucks selbstständig aufgebaut.}
					\item Erkundung einer unbekannten rechtwinkligen Fläche ausgehend von fest definierten Startpositionen
						\\ \textsl{Das Spielfeld besteht aus quadratischen Feldern. Diese sind rasterförmig und
						   zusammenhängend angeordnet. Die Felder sind mindestens so groß, dass ein im Zentrum des Quadrats
						   stehender e-puck vollständig in das Innere des Feldes passt. Die Startpositionen der e-puck Roboter sind
						   fest vordefiniert.}
					\item Fortbewegung auf den Kanten der Quadrate
						\\ \textsl{Die Linien des Spielfeldes müssen schwarz mit ausreichendem Kontrastverhältnis
						    zum Untergrund sein. Zwingend erforderlich ist, dass die Breite der Linien innerhalb der Spezifikation der
							Bodensensoren liegt.}		
					\item Vermeidung von Kollisionen
						\\ \textsl{Während der Bewegung über das Spielfeld vermeiden die Roboter Kollisionen mit anderen Robotern.}	
					\item Gleichberechtigung der teilnehmenden Roboter
						\\ \textsl{Alle Roboter der Gruppe haben den selben Aufgabenbereich. Es gibt keine zentrale Einheit
							zur Koordination und Synchronisation.}	
					\item Rückkehr zum Ausgangspunkt
						\\ \textsl{Nach Abschluss des Erkundungsvorgangs kehren alle Roboter zu ihren jeweiligen Startpositionen
							zurück.}	
				\end{itemize}
				\item Smartphone
				\begin{itemize}
					\item Kommunikation über Bluetooth
						\\ \textsl{Die Verbindung und Kommunikation mit den Robotern erfolgt über die Bluetooth-Schnittstelle des
							Smartphones.}
					\item Visualisierung
						\\ \textsl{Die bereits erkundeten Gebiete werden in einer benutzerfreundlichen Android-Anwendung übersichtlich
							dargestellt. Auf dieser Karte werden zusätzlich die aktuellen Positionen der e-pucks eingetragen.}
					\item Manuelle Regelung der Bewegungsgeschwindigkeit
						\\ \textsl{Mit dem Smartphone kann die Geschwindigkeit der teilnehmenden, automatisch erkundenden e-pucks stufenweise 										eingestellt werden.}	
					\item Manuelle Steuerung eines Roboters
						\\ \textsl{Mit Hilfe der Anwendung kann einer der teilnehmenden e-pucks ausgewählt werden. Dieser lässt sich
							durch den Benutzer entlang der Kanten steuern. Es kann jeweils nur ein Roboter zur selben Zeit ausgewählt
							und gesteuert werden.}		
					\item Zwei Steuerungsarten
						\\ \textsl{Die manuelle Steuerung der Roboter erfolgt wahlweise über einen On-Screen-Joystick oder über
							den im Handy integrierten Beschleunigungssensor.}											
				\end{itemize}
			\end{itemize}
		\subsection{Wunschkriterien}
			\begin{itemize}
				\item e-puck Roboter
				\begin{itemize}
					\item Kritische Bereiche
						\\ \textsl{Auf dem Spielfeld gibt es Bereiche, in denen ein e-puck keine Nachrichten senden
							und empfangen darf. Diese Gebiete werden als `Kritische Bereiche' bezeichnet und müssen
							während des Erkundungsvorgangs besonders berücksichtigt werden.}											
					\item Beliebige Startpositionen
						\\ \textsl{Die Roboter können auf frei wählbaren Startpositionen innerhalb des Spielfeldes
							abgesetzt werden. Um dennoch Synchronisation zu erreichen ist ein erweiterter Lokalisierungsvorgang
							notwendig.}
					\item Zustandsvisualisierung
						\\ \textsl{Die e-pucks stellen ihren aktuellen Zustand, zum Beispiel `Erkundung läuft' oder `Erkundung
							beendet', mit Hilfe der ihnen zur Verfügung stehenden Mittel optisch dar.}												
					\item Einstellbarer Synchronisationsmodus
						\\ \textsl{Wirklich sinnvoll?!}		
				\end{itemize}
				\item Smartphone
				\begin{itemize}
					\item Anzeige von Statusinformationen
						\\ \textsl{Mit Hilfe des Smartphones ist es möglich, Informationen über den aktuellen Status eines e-pucks
							übersichtlich anzuzeigen.}		
					\item Pfadanzeige von einzelnen Robotern
						\\ \textsl{Das Smartphone bietet die Möglichkeit, den Pfad anzuzeigen, den jeder e-puck im Rahmen der
							Erkundung zurückgelegt hat.}							
					\item Exportfunktion für erkundete Karten
						\\ \textsl{Teilweise und vollständig erkundete Karten können in geeigneter Weise exportiert und auf dem
							Smartphone abgespeichert werden.}		
					\item Internationalisierung
						\\ \textsl{Die Anwendung ist in den Sprache Deutsch und Englisch verfügbar.}									
				\end{itemize}
			\end{itemize}
		\subsection{Abgrenzungskriterien}
			\begin{itemize}
				\item Keine Unterstützung für abweichende Spielfelder
					\\ \textsl{Die Felder und zugehörigen Linien müssen den in den Musskriterien genannten Voraussetzungen
						genügen. Es sind zum Beispiel keine runden oder diagonalen Verbindungslinien erlaubt.}
				\item Keine Berücksichtigung von dynamischen Änderungen des Spielfeldes
					\\ \textsl{Nachdem der Erkundungsvorgang gestartet wurde, dürfen keine Modifikationen am
						Spielfeld getätigt werden, welche Einfluss auf den Erkundungsalgorithmus zur Folge hätten.
						Insbesondere werden keine nachträglichen Änderungen an bereits erkundeten Gebieten berücksichtigt.}			
				\item Unterstützung für maximal ein Smartphone
					\\ \textsl{Es kann lediglich ein Smartphone zur Auswahl, Steuerung und Visualisierung verwendet
						werden.}		
			\end{itemize}
	\section{Produkteinsatz}
		\subsection{Anwendungsbereiche}
			Das Projekt ist eine Forschungsarbeit in den Gebieten:
			\begin{itemize}
				\item Robotik
				\item Verteilte Systeme
				\item Künstliche Intelligenz
			\end{itemize}
		\subsection{Zielgruppen}
			Zielgruppe des Projektes sind insbesondere: 
			\begin{itemize}
				\item Studenten
				\item Forschungsgruppen in ähnlichen Bereichen
				\item e-puck Community
			\end{itemize}
		\subsection{Betriebsbedingungen}
			\begin{itemize}
				\item Ausreichende Stromversorgung
					\\ \textsl{Die Akkuleistung der einzelnen Roboter und des Smartphones muss für die gesamte
						Dauer der Lokalisierung, Erkundung und Rückkehr in die Startpositionen ausreichend sein.}
				\item Geeignete Bedingungen für Funknetzwerke
					\\ \textsl{Die Ausmaße des Spielfelds, die Abstände zwischen den e-pucks und zum Smartphone sowie
						Signale anderer Netze dürfen keine störenden Einflüsse auf die Bluetooth-Verbindungen haben.}
				\item Größe des Spielfeldes
					\\ \textsl{Der integrierte Arbeitsspeicher des e-puck Roboter stellt eine Begrenzung des 
					 persistierbaren lokalen Spielfeldes dar. Es werden ca. 700 Knotenpunkte unterstützt. 
					 !!! Diese Rechnung ist sehr fragwürdig !!!}
				\item Betriebsbedingungen der e-puck Roboter und des Smartphones
					\\ \textsl{Die weiteren Betriebsbedingungen können dem Benutzerhandbuch entnommen werden.}			
				\item Wartungsfrei
					\\ \textsl{Das System bedarf keiner regelmäßigen Wartung oder Aktualisierung.}		
			\end{itemize}
<<<<<<< .mine
	\section{Produktumgebung}
		\subsection{Software}
			\begin{itemize}
				\item Geeigneter Bootloader oder Programmer zum Programmieren der e-puck Roboter
				\item Android Software ab Version 2.1
			\end{itemize}
		\subsection{Hardware}
			\begin{itemize}
				\item e-puck Roboter
				\item Bodensensor für e-puck Roboter
				\item Android Smartphone
					\\ \textsl{Voraussetzungen: Bluetooth, geeignete Auflösung, Beschleunigungssensoren, Touch-Display.}
				\item Computer mit Bluetooth Unterstützung
			\end{itemize}		
	\section{Produktfunktionen}
		\subsection{e-puck Roboter}
			\subsubsection{Netzwerkfunktionen}
				\begin{list}{\labelitemi}{\leftmargin=1cm}
					\item[\textbf{\textbackslash F50\textbackslash}] Suchfunktion
						\\ \textsl{Ein e-puck Roboter kann andere e-puck Roboter in Reichweite seines Bluetooth-Moduls entdecken.}
					\item[\textbf{\textbackslash F60\textbackslash}] Verbindungsaufbau
						\\ \textsl{Ein e-puck Roboter kann zu einem anderen e-puck Roboter, den er mit der Suchfunktion entdeckt
						hat, eine Bluetooth-Verbindung aufbauen.}
					\item[\textbf{\textbackslash F70\textbackslash}] Broadcast-Kommunikation
						\\ \textsl{Die e-puck Roboter können über bestehende Bluetooth-Verbindungen per Broadcast
						Nachrichten austauschen. Auch die Kommunikation mit dem Smartphone erfolgt via Broadcast.}
					\item[\textbf{\textbackslash F75W\textbackslash}] Kommunikationsunterdrückung
						\\ \textsl{Erreicht ein e-puck einen kritischen Bereich auf dem Spielfeld, kann er über bestehende
						Bluetooth-Verbindungen weder Nachrichten senden noch empfangen. Die Verbindungen an sich bleiben bestehen.}
				\end{list}
			\subsubsection{Bewegungsfunktionen}
				\begin{list}{\labelitemi}{\leftmargin=1cm}
					\item[\textbf{\textbackslash F80\textbackslash}] Fahrsteuerung
						\\ \textsl{}
					\item[\textbf{\textbackslash F100\textbackslash}] Linienverfolgung
						\\ \textsl{Die e-pucks bewegen sich auf dem Spielfeld nur auf den vorhandenen Linien. Das gilt auch für e-pucks,
						die mit dem Smartphone halbmanuell gesteuert werden.}
					\item[\textbf{\textbackslash F90\textbackslash}] Kollisionsvermeidung
						\\ \textsl{Ein e-puck erkennt und vermeidet direkt bevorstehende Kollisionen mit anderen e-pucks.}
					\item[\textbf{\textbackslash F110\textbackslash}] Knotenanalyse
						\\ \textsl{Ein e-puck erkennt beim Überfahren eines Knotens, in welche Richtungen er weiterfahren kann.}
				\end{list}
			\subsubsection{Explorationsfunktionen}
				\begin{list}{\labelitemi}{\leftmargin=1cm}
					\item[\textbf{\textbackslash F120\textbackslash}] Lokalisierung
						\\ \textsl{Vor Beginn des Erkundungsvorgangs einigen sich die e-pucks untereinander auf einen gemeinsamen Bezugspunkt
						und ermitteln relativ dazu ihre eigene Position. Die e-pucks starten in diesem Fall von fest definierten Startpositionen.}
					\item[\textbf{\textbackslash F150\textbackslash}] Erkundungsfunktion
						\\ \textsl{Nach Abschluss des Lokalisierungsvorgangs erkunden die e-pucks in Zusammenarbeit das unbekannte
							Spielfeld vollständig.}
					\item[\textbf{\textbackslash F140\textbackslash}] Kartensynchronisation
						\\ \textsl{Während des Erkundungsvorgangs tauschen die e-pucks kontinuierlich Informationen über die von ihnen
						erkundeten Felder aus. Ist ein Smartphone mit dem Netzwerk verbunden, erhält auch dieses entsprechende Nachrichten.}
					\item[\textbf{\textbackslash F130\textbackslash}] Lokale Kartenkonstruktion	
						\\ \textsl{Jeder e-puck speichert lokal die komplette bisher erkundete Karte. Diese wird mit Hilfe
						der Informationen aus den Nachrichten der anderen e-pucks sowie den eigenen Erkundungsergebnissen konstruiert.} 				
					\item[\textbf{\textbackslash F160\textbackslash}] Rückkehrfunktion
						\\ \textsl{Nachdem das Gebiet vollständig erkundet worden ist, fahren die e-pucks wieder an ihre Startpositionen
						zurück.}
					\item[\textbf{\textbackslash F170W\textbackslash}] Globale Lokalisierung (evtl. inkl. Steuerung)
						\\ \textsl{Starten die e-pucks nicht auf fest definierten Startpositionen, sondern auf beliebigen Feldern innerhalb
							des Spielfeldes, sammeln sich die Roboter zunächst nebeneinander in einem Gebiet des Spielfeldes. Dort beginnt der
							Erkundungsvorgang mit der oben beschriebenen Lokalisierung.}
				\end{list}
		\subsection{Smartphone}
			\subsubsection{Netzwerkfunktionen}
				\begin{list}{\labelitemi}{\leftmargin=1cm}
					\item[\textbf{\textbackslash F200\textbackslash}] Suchfunktion
						\\ \textsl{Das Smartphone kann e-puck Roboter in Reichweite seines Bluetooth-Moduls entdecken.}
					\item[\textbf{\textbackslash F210\textbackslash}] Verbindungsaufbau
						\\ \textsl{Das Smartphone kann sich via Bluetooth mit einem e-puck verbinden, den es mit Hilfe
						der Suchfunktion entdeckt hat.}				
					\item[\textbf{\textbackslash F220\textbackslash}] Kommunikationsfunktion
						\\ \textsl{Das Smartphone kann Daten über eine korrekte Bluetoothverbindung mit einem e-puck Roboter Daten senden und
							empfangen.}
				\end{list}
			\subsubsection{Steuerungsfunktionen}		
				\begin{list}{\labelitemi}{\leftmargin=1cm}
					\item[\textbf{\textbackslash F290\textbackslash}] Steuerungsfunktion mit einem On-Screen-Joystick
						\\ \textsl{Der ausgewählte e-puck kann mit einem On-Screen-Joystick halbmanuell gesteuert werden. Die Bewegung
							erfolgt weiterhin auf den Linien des Spielfeldes.}	
					\item[\textbf{\textbackslash F300\textbackslash}] Steuerungsfunktion mit dem eingebauten Beschleunigungssensoren
						\\ \textsl{Alternativ kann der ausgewählte e-puck kann mit Hilfe des eingebauten Beschleunigungssensors halbmanuell
						gesteuert werden. Auch in diesem Fall bewegt sich der Roboter nur auf den Linien des Spielfeldes.}					
				\end{list}					
			\subsubsection{Grafische Benutzeroberfläche}
				\begin{list}{\labelitemi}{\leftmargin=1cm}
					\item[\textbf{\textbackslash F230\textbackslash}] Kartendarstellung
						\\ \textsl{Die GUI stellt die bereits erkundeten Gebiete in geeigneter Art und Weise dar.}	
					\item[\textbf{\textbackslash F240\textbackslash}] Roboterdarstellung
						\\ \textsl{Die aktuelle Position der e-pucks wird auf der Karte der bisher erkundeten Gebiete
							einzeichnet.}	
					\item[\textbf{\textbackslash F250\textbackslash}] Steuerungsfunktion per On-Screen-Joystick
						\\ \textsl{Mithilfe eines Steuerknüppels kann der Teilnehmer entlang der Linien bewegt werden}
					\item[\textbf{\textbackslash F260\textbackslash}] Steuerungsfunktion per Bewegungssensor	
						\\ \textsl{Durch Kippen des Smartphones kann der e-puck Roboter beschleunigt, gebremst sowie in verschiedene, gültige
							Richtungen bewegt werden.}
					\item[\textbf{\textbackslash F270W\textbackslash}] Zustandsvisualisierung (status + pfade)		
						\\ \textsl{Auf dem Smartphone wird für jeden e-puck Roboter der bisher abgefahrene Weg, die Anzahl der überfahreren
							Knoten und sein momentaner Status (stehend, fahrend, manuell gesteuert) visualisiert.}
					\item[\textbf{\textbackslash F280W\textbackslash}] Internationalisierung	
						\\ \textsl{Es gibt die Möglichkeit zwischen deutscher und englischer Sprachausgabe zu wählen.}
				\end{list}
				
	\section{Produktdaten}
	
	
	\section{Produktleistungen}
	
	
	\section{Benutzungsoberfläche}
	
	
	\section{Qualitätsbestimmungen}
	
	
	\section{Globale Testszenarien und Testfälle}
	
	
	\section{Produktfunktionen}	
	
	
	\section{Ergänzungen}
						
=======
			
			
>>>>>>> .r8
\end{document}